\documentclass[10pt,a4paper]{amsart}
\usepackage{amsmath,amssymb,amsthm,graphicx}
% Retained preamble for the Stacks Project chapter snapshot in this directory.
% The chapter's own preamble.tex is a copy of the shared Stacks Project preamble; the block
% below is sliced MECHANICALLY, line for line, from it. Nothing was added, removed or
% redefined; the wrapper class, the xr-hyper/hyperref wiring and the externaldocument list
% are omitted because the assembler supplies the document class.
%
% lines taken from preamble.tex: 3, 5, 9, 11, 12, 156, 157, 158, ... (51 lines total)
%
% package coverage restored from the chapter's own preamble (2026-10-13): xy, xy-2cell

\usepackage[all]{xy}\usepackage{verbatim}
\xyoption{2cell}
\UseAllTwocells

\newenvironment{reference}{\comment}{\endcomment}
\newenvironment{slogan}{\comment}{\endcomment}
\newenvironment{history}{\comment}{\endcomment}
\theoremstyle{plain}
\newtheorem{theorem}[subsection]{Theorem}
\newtheorem{proposition}[subsection]{Proposition}
\newtheorem{lemma}[subsection]{Lemma}
\theoremstyle{definition}
\newtheorem{definition}[subsection]{Definition}
\newtheorem{example}[subsection]{Example}
\newtheorem{exercise}[subsection]{Exercise}
\newtheorem{situation}[subsection]{Situation}
\theoremstyle{remark}
\newtheorem{remark}[subsection]{Remark}
\newtheorem{remarks}[subsection]{Remarks}
\numberwithin{equation}{subsection}
\def\lim{\mathop{\mathrm{lim}}\nolimits}
\def\colim{\mathop{\mathrm{colim}}\nolimits}
\def\Spec{\mathop{\mathrm{Spec}}}
\def\Hom{\mathop{\mathrm{Hom}}\nolimits}
\def\Ext{\mathop{\mathrm{Ext}}\nolimits}
\def\SheafHom{\mathop{\mathcal{H}\!\mathit{om}}\nolimits}
\def\SheafExt{\mathop{\mathcal{E}\!\mathit{xt}}\nolimits}
\def\Sch{\mathit{Sch}}
\def\Mor{\mathop{\mathrm{Mor}}\nolimits}
\def\Ob{\mathop{\mathrm{Ob}}\nolimits}
\def\Sh{\mathop{\mathit{Sh}}\nolimits}
\def\NL{\mathop{N\!L}\nolimits}
\def\CH{\mathop{\mathrm{CH}}\nolimits}
\def\proetale{{pro\text{-}\acute{e}tale}}
\def\etale{{\acute{e}tale}}
\def\QCoh{\mathit{QCoh}}
\def\Ker{\mathop{\mathrm{Ker}}}
\def\Im{\mathop{\mathrm{Im}}}
\def\Coker{\mathop{\mathrm{Coker}}}
\def\Coim{\mathop{\mathrm{Coim}}}
\DeclareMathSymbol{\boxtimes}{\mathbin}{AMSa}{"02}
\def\QCohstack{\mathcal{QC}\!\mathit{oh}}
\def\Cohstack{\mathcal{C}\!\mathit{oh}}
\def\Spacesstack{\mathcal{S}\!\mathit{paces}}
\def\Quotfunctor{\mathrm{Quot}}
\def\Hilbfunctor{\mathrm{Hilb}}
\def\Curvesstack{\mathcal{C}\!\mathit{urves}}
\def\Polarizedstack{\mathcal{P}\!\mathit{olarized}}
\def\Complexesstack{\mathcal{C}\!\mathit{omplexes}}
\def\Pic{\mathop{\mathrm{Pic}}\nolimits}
\def\Picardstack{\mathcal{P}\!\mathit{ic}}
\def\Picardfunctor{\mathrm{Pic}}
\def\Deformationcategory{\mathcal{D}\!\mathit{ef}}
\usepackage{cleveref}
\crefname{originallemma}{original lemma}{original lemmas}
\Crefname{originallemma}{Original lemma}{Original lemmas}
\makeatletter
\newlabel{derived-lemma-adjoint-preserve-K-injectives}{{08BJ}{source}{}{sourcepointer-1}{}}
\newlabel{derived-lemma-adjoint-preserve-K-injectives@cref}{{[originallemma][0][]08BJ}{[1][0][]source}{}{}{}}
\makeatother
\crefname{originallemma}{original lemma}{original lemmas}
\Crefname{originallemma}{Original lemma}{Original lemmas}
\makeatletter
\newlabel{lemma-affine-pushforward}{{08GY}{source}{}{sourcepointer-2}{}}
\newlabel{lemma-affine-pushforward@cref}{{[originallemma][0][]08GY}{[1][0][]source}{}{}{}}
\makeatother
\crefname{originallemma}{original lemma}{original lemmas}
\Crefname{originallemma}{Original lemma}{Original lemmas}
\makeatletter
\newlabel{lemma-induction-principle-separated}{{08GQ}{source}{}{sourcepointer-3}{}}
\newlabel{lemma-induction-principle-separated@cref}{{[originallemma][0][]08GQ}{[1][0][]source}{}{}{}}
\makeatother
\crefname{originallemma}{original lemma}{original lemmas}
\Crefname{originallemma}{Original lemma}{Original lemmas}
\makeatletter
\newlabel{proposition-quasi-compact-affine-diagonal}{{08H1}{source}{}{sourcepointer-4}{}}
\newlabel{proposition-quasi-compact-affine-diagonal@cref}{{[originallemma][0][]08H1}{[1][0][]source}{}{}{}}
\makeatother
\crefname{originallemma}{original lemma}{original lemmas}
\Crefname{originallemma}{Original lemma}{Original lemmas}
\makeatletter
\newlabel{remark-how-to}{{08GR}{source}{}{sourcepointer-5}{}}
\newlabel{remark-how-to@cref}{{[originallemma][0][]08GR}{[1][0][]source}{}{}{}}
\makeatother
\crefname{originallemma}{original lemma}{original lemmas}
\Crefname{originallemma}{Original lemma}{Original lemmas}
\makeatletter
\newlabel{spaces-morphisms-lemma-affine-permanence}{{08GB}{source}{}{sourcepointer-6}{}}
\newlabel{spaces-morphisms-lemma-affine-permanence@cref}{{[originallemma][0][]08GB}{[1][0][]source}{}{}{}}
\makeatother
\crefname{originallemma}{original lemma}{original lemmas}
\Crefname{originallemma}{Original lemma}{Original lemmas}
\makeatletter
\newlabel{spaces-properties-definition-separated}{{03BS}{source}{}{sourcepointer-7}{}}
\newlabel{spaces-properties-definition-separated@cref}{{[originallemma][0][]03BS}{[1][0][]source}{}{}{}}
\makeatother
\title{Stacks, Derived Categories of Spaces: direct-image-coherator}
\author{The Stacks Project Authors}
\date{stacks-project:spaces-perfect.tex:local-snapshot-20261012; GNU FDL 1.2 or later}
\begin{document}
\maketitle
\noindent\textit{Editorial note.} Statement span is lines 2713-2734 of sources/stacks-spaces-perfect/pinned-original-source.tex; proof span is lines 2736-2813, the complete author proof (76 body lines).\par\medskip
\noindent\textit{Editorial note.} Navigation stubs: 7 cited-but-not-included labels are supplied as non-load-bearing pointers with their pinned tags where the snapshot tag index has them.\par\medskip
\noindent\textit{Editorial note.} extra\_packages is empty on purpose: the retained preamble already loads xy with options and a second declaration clashes in the assembler wrapper.\par\medskip
\noindent\textit{Editorial note.} Text is retained verbatim, including the author's own typos and placeholders.\par\medskip
\noindent\textit{Editorial note.} Queue index 18 of queue\_spaceperfect.json (pb=3409, stmt 2713-2734, proof 2736-2813).\par\medskip
\noindent\textit{Editorial note.} stacks-spaces-perfect wave sp01 (2026-10-14): fresh adjudication accept per lane-s3/sp01\_judgments; pool per Chief 12:4x unfreeze (reserve-pool workflow, six-check clean); spaces-perfect.tex 126 with-proof rows, 70 b107-occupied labels avoided, 56 free rows judged 32 accept / 24 reject.\par\medskip
\noindent\textit{Editorial note.} Cross-chapter same-label warning (NOT an occupancy conflict, Chief ruling 2026-10-13 evening): label 'lemma-direct-image-coherator' also exists in stacks-perfect — a different proposition in a different chapter file.\par\medskip
\noindent\textit{Allowed standard prerequisites (editorial).}\begin{itemize}
\item Explicit prerequisite boundary. The cited inputs are prerequisites; the argument in the retained span is not replaced by any cited result.
\end{itemize}\medskip
\ifdefined\hypertarget\hypertarget{sourcepointer-1}{}\fi
\noindent\textit{Source reference.} Original lemma 08BJ: sources/stacks-spaces-perfect/pinned-original-source.tex, source lines 2813--2813. Cited by the retained statement or proof but not part of the collected spans; supplied only as a navigation pointer. The argument does not depend on its content.\par\medskip
\ifdefined\hypertarget\hypertarget{sourcepointer-2}{}\fi
\noindent\textit{Source reference.} Original lemma 08GY: sources/stacks-spaces-perfect/pinned-original-source.tex, source lines 2339--2353. Cited by the retained statement or proof but not part of the collected spans; supplied only as a navigation pointer. The argument does not depend on its content.\par\medskip
\ifdefined\hypertarget\hypertarget{sourcepointer-3}{}\fi
\noindent\textit{Source reference.} Original lemma 08GQ: sources/stacks-spaces-perfect/pinned-original-source.tex, source lines 1643--1662. Cited by the retained statement or proof but not part of the collected spans; supplied only as a navigation pointer. The argument does not depend on its content.\par\medskip
\ifdefined\hypertarget\hypertarget{sourcepointer-4}{}\fi
\noindent\textit{Source reference.} Original lemma 08H1: sources/stacks-spaces-perfect/pinned-original-source.tex, source lines 2689--2700. Cited by the retained statement or proof but not part of the collected spans; supplied only as a navigation pointer. The argument does not depend on its content.\par\medskip
\ifdefined\hypertarget\hypertarget{sourcepointer-5}{}\fi
\noindent\textit{Source reference.} Original lemma 08GR: sources/stacks-spaces-perfect/pinned-original-source.tex, source lines 1677--1700. Cited by the retained statement or proof but not part of the collected spans; supplied only as a navigation pointer. The argument does not depend on its content.\par\medskip
\ifdefined\hypertarget\hypertarget{sourcepointer-6}{}\fi
\noindent\textit{Source reference.} Original lemma 08GB: sources/stacks-spaces-perfect/pinned-original-source.tex, source lines 2813--2813. Cited by the retained statement or proof but not part of the collected spans; supplied only as a navigation pointer. The argument does not depend on its content.\par\medskip
\ifdefined\hypertarget\hypertarget{sourcepointer-7}{}\fi
\noindent\textit{Source reference.} Original lemma 03BS: sources/stacks-spaces-perfect/pinned-original-source.tex, source lines 2813--2813. Cited by the retained statement or proof but not part of the collected spans; supplied only as a navigation pointer. The argument does not depend on its content.\par\medskip
\begin{lemma}
\label{lemma-direct-image-coherator}
Let $S$ be a scheme and let $f : X \to Y$ be a morphism of algebraic
spaces over $S$. Assume $X$ and $Y$ are quasi-compact and have affine diagonal
over $\mathbf{Z}$ (as in Properties of Spaces, Definition
\ref{spaces-properties-definition-separated}). Then, denoting
$$
\Phi : D(\QCoh(\mathcal{O}_X)) \to D(\QCoh(\mathcal{O}_Y))
$$
the right derived functor of
$f_* : \QCoh(\mathcal{O}_X) \to \QCoh(\mathcal{O}_Y)$
the diagram
$$
\xymatrix{
D(\QCoh(\mathcal{O}_X)) \ar[d]_\Phi \ar[r] &
D_\QCoh(\mathcal{O}_X) \ar[d]^{Rf_*} \\
D(\QCoh(\mathcal{O}_Y)) \ar[r] &
D_\QCoh(\mathcal{O}_Y)
}
$$
is commutative.
\end{lemma}

\begin{proof}
Observe that the horizontal arrows in the diagram are
equivalences of categories by
Proposition \ref{proposition-quasi-compact-affine-diagonal}.
Hence we can identify these categories (and similarly for
other quasi-compact algebraic spaces with affine diagonal)
and then the statement of the lemma is that the canonical map
$\Phi(K) \to Rf_*(K)$ is an isomorphism for all $K$ in
$D(\QCoh(\mathcal{O}_X))$. Note that if $K_1 \to K_2 \to K_3 \to K_1[1]$
is a distinguished triangle in $D(\QCoh(\mathcal{O}_X))$ and
the statement is true for two-out-of-three, then it is true
for the third.

\medskip\noindent
Let $\mathcal{B} \subset \Ob(X_{spaces, \etale})$ be the set of
objects which are quasi-compact and have affine diagonal.
For $U \in \mathcal{B}$ and any morphism $g : U \to Z$
where $Z$ is a quasi-compact algebraic space over $S$ with
affine diagonal, denote
$$
\Phi_g : D(\QCoh(\mathcal{O}_U)) \to D(\QCoh(\mathcal{O}_Z))
$$
the derived extension of $g_*$. Let
$P(U) =$ ``for any $K$ in $D(\QCoh(\mathcal{O}_U))$
and any $g : U \to Z$ as above the map $\Phi_g(K) \to Rg_*K$
is an isomorphism''.
By Remark \ref{remark-how-to} conditions (1), (2), and (3)(a) of
Lemma \ref{lemma-induction-principle-separated} hold and we are
left with proving (3)(b) and (4).

\medskip\noindent
Checking condition (3)(b). Let $U$ be an affine scheme \'etale
over $X$. Let $g : U \to Z$ be as above. Since the diagonal of $Z$
is affine the morphism $g : U \to Z$
is affine (Morphisms of Spaces, Lemma
\ref{spaces-morphisms-lemma-affine-permanence}).
Hence $P(U)$ holds by Lemma \ref{lemma-affine-pushforward}.

\medskip\noindent
Checking condition (4).
Let $(U \subset W, V \to W)$ be an elementary distinguished square
in $X_{spaces, \etale}$ with $U, W, V$ in $\mathcal{B}$ and $V$ affine.
Assume that $P$ holds for $U$, $V$, and $U \times_W V$.
We have to show that $P$ holds for $W$. Let $g : W \to Z$
be a morphism to a quasi-compact algebraic space with affine diagonal.
Let $K$ be an object of $D(\QCoh(\mathcal{O}_W))$.
Consider the distinguished triangle
$$
K \to Rj_{U, *}K|_U \oplus Rj_{V, *}K|_V \to
Rj_{U \times_W V, *}K|_{U \times_W V} \to K[1]
$$
in $D(\mathcal{O}_W)$. By the two-out-of-three property mentioned
above, it suffices to show that $\Phi_g(Rj_{U, *}K|_U) \to Rg_*(Rj_{U, *}K|_U)$
is an isomorphism and similarly for $V$ and $U \times_W V$.
This is discussed in the next paragraph.

\medskip\noindent
Let $j : U \to W$ be a morphism $X_{spaces, \etale}$ with
$U, W$ in $\mathcal{B}$ and $P$ holds for $U$. Let $g : W \to Z$
be a morphism to a quasi-compact algebraic space with affine diagonal.
To finish the proof we have to show that
$\Phi_g(Rj_*K) \to Rg_*(Rj_*K)$
is an isomorphism for any $K$ in $D(\QCoh(\mathcal{O}_U))$.
Let $\mathcal{I}^\bullet$ be a K-injective complex in $\QCoh(\mathcal{O}_U)$
representing $K$.
From $P(U)$ applied to $j$ we see that
$j_*\mathcal{I}^\bullet$ represents $Rj_*K$.
Since $j_* : \QCoh(\mathcal{O}_U) \to \QCoh(\mathcal{O}_X)$
has an exact left adjoint
$j^* : \QCoh(\mathcal{O}_X) \to \QCoh(\mathcal{O}_U)$
we see that $j_*\mathcal{I}^\bullet$ is a K-injective complex
in $\QCoh(\mathcal{O}_W)$, see
Derived Categories, Lemma \ref{derived-lemma-adjoint-preserve-K-injectives}.
Hence $\Phi_g(Rj_*K)$ is represented by
$g_*j_*\mathcal{I}^\bullet = (g \circ j)_*\mathcal{I}^\bullet$.
By $P(U)$ applied to $g \circ j$ we see that this represents
$R_{g \circ j, *}(K) = Rg_*(Rj_*K)$. This finishes the proof.
\end{proof}

\end{document}
